---
tipo: Discursiva
dificuldade: Média
nivel: médio
disciplina: Matemática
assunto: Trigonometria
gabarito:
resposta: x = 2k\pi - \frac{\pi}{2},\ k \in \mathbb{Z}
tags: [desigualdade, cosseno, seno]
fonte: concurso
concurso:
banca: UFF
ano: 1999
numero: 1
cargo:
prova:
---
\question
Determine o(s) valor(es) de $x \in \mathbb{R}$ que satisfaz(em) à desigualdade:

$\cos^2 x \ge 2(\sin x + 1)$

---
tipo: Discursiva
dificuldade: Média
nivel: médio
disciplina: Matemática
assunto: Progressão Geométrica
gabarito:
resposta: x = 2 ou x = 0
tags: [série geométrica]
fonte: concurso
concurso:
banca: UFF
ano: 1999
numero: 2
cargo:
prova:
---
\question
Considere

$S = (x - 1)^2 + \frac{(x - 1)^2}{2} + \frac{(x - 1)^2}{4} + \frac{(x - 1)^2}{8} + \cdots$

Determine o(s) valor(es) de $x$ que torna(m) $S = 2$.

---
tipo: Múltipla Escolha
dificuldade: Fácil
nivel: médio
disciplina: Matemática
assunto: Números Inteiros
gabarito: D
resposta:
tags: [paridade]
fonte: concurso
concurso:
banca: UFF
ano: 2000
numero: 3
cargo:
prova:
---
\question
Considere $p$, $q \in \mathbb{N}^*$ tais que $p$ e $q$ são números pares. Se $p > q$, pode-se afirmar que:
\begin{choices}
\choice $(pq + 1)$ é múltiplo de 4.
\choice $p - q$ é ímpar.
\choice $p + q$ é primo.
\CorrectChoice $p^2 - q^2$ é par.
\choice $p(q + 1)$ é ímpar.
\end{choices}

---
tipo: Múltipla Escolha
dificuldade: Média
nivel: médio
disciplina: Matemática
assunto: Geometria Analítica
gabarito: B
resposta:
tags: [equação da reta, figura]
fonte: concurso
concurso:
banca: UFF
ano: 2000
numero: 4
cargo:
prova:
---
\question
Na figura a seguir estão representadas as retas $r$ e $s$.

\includegraphics{figura-1}

Sabendo que a equação da reta $s$ é $x = 3$ e que $OP$ mede 5 cm, a equação de $r$ é:
\begin{choices}
\choice $y = \frac{3x}{4}$
\CorrectChoice $y = \frac{4x}{3}$
\choice $y = \frac{5x}{3}$
\choice $y = 3x$
\choice $y = 5x$
\end{choices}

---
tipo: Múltipla Escolha
dificuldade: Média
nivel: médio
disciplina: Matemática
assunto: Teoria dos Conjuntos
gabarito: B
resposta:
tags: [inteiros, inclusão]
fonte: concurso
concurso:
banca: UFF
ano: 2000
numero: 5
cargo:
prova:
---
\question
Com relação aos conjuntos

$P = \{x \in \mathbb{Z} \mid |x| \le \sqrt{7}\}$ e

$Q = \{x \in \mathbb{Z} \mid x^2 \le 0{,}333\ldots\}$

afirma-se:

I. $P \cup Q = P$

II. $Q - P = \{0\}$

III. $P \subset Q$

IV. $P \cap Q = Q$

Somente são verdadeiras as afirmativas:
\begin{choices}
\choice I e III
\CorrectChoice I e IV
\choice II e III
\choice II e IV
\choice III e IV
\end{choices}

---
tipo: Múltipla Escolha
dificuldade: Média
nivel: médio
disciplina: Matemática
assunto: Geometria Plana
gabarito: E
resposta:
tags: [área, retângulo, paralelogramo, figura]
fonte: concurso
concurso:
banca: UFF
ano: 2000
numero: 6
cargo:
prova:
---
\question
Na figura, $MNPQ$ é um retângulo, $MNUV$ é um paralelogramo, as medidas de $MQ$ e $MV$ são iguais e $0^\circ < \alpha < 45^\circ$.

\includegraphics{figura-2}

Indicando-se por $S$ a área de $MNPQ$ e por $S'$ a área de $MNUV$, conclui-se que:
\begin{choices}
\choice $S = S' \sen \alpha$
\choice $S' = S$
\choice $S' = S \cos \alpha$
\choice $S = S' \cos \alpha$
\CorrectChoice $S' = S \sen \alpha$
\end{choices}

---
tipo: Múltipla Escolha
dificuldade: Fácil
nivel: médio
disciplina: Matemática
assunto: Média Aritmética
gabarito: A
resposta:
tags: [média]
fonte: concurso
concurso:
banca: UFF
ano: 2000
numero: 7
cargo:
prova:
---
\question
Para que a média aritmética das notas de uma turma de 20 alunos aumentasse em 0,1, alterou-se uma dessas notas para 7,5. Antes da alteração, tal nota era:
\begin{choices}
\CorrectChoice 5,5
\choice 6,0
\choice 7,4
\choice 7,6
\choice 8,5
\end{choices}

---
tipo: Múltipla Escolha
dificuldade: Média
nivel: médio
disciplina: Matemática
assunto: Trigonometria
gabarito: E
resposta:
tags: [círculo trigonométrico, figura]
fonte: concurso
concurso:
banca: UFF
ano: 2000
numero: 8
cargo:
prova:
---
\question
Considere os ângulos $\alpha$, $\beta$ e $\gamma$ conforme representados no círculo.

\includegraphics{figura-3}

Pode-se afirmar que:
\begin{choices}
\choice $\cos \alpha < \cos \beta$
\choice $\cos \gamma > \cos \alpha$
\choice $\sen \alpha > \sen \beta$
\choice $\sen \beta < \cos \gamma$
\CorrectChoice $\cos \beta < \cos \gamma$
\end{choices}

---
tipo: Múltipla Escolha
dificuldade: Média
nivel: médio
disciplina: Matemática
assunto: Geometria Espacial
gabarito: D
resposta:
tags: [tetraedro regular, segmento médio, figura]
fonte: concurso
concurso:
banca: UFF
ano: 2000
numero: 9
cargo:
prova:
---
\question
No tetraedro regular representado na figura, $R$ e $S$ são, respectivamente, os pontos médios de $NP$ e $OM$.

\includegraphics{figura-4}

A razão $\frac{RS}{MN}$ é igual a:
\begin{choices}
\choice $\sqrt{3}$
\choice $\frac{\sqrt{3}}{2}$
\choice $\sqrt{2}$
\CorrectChoice $\frac{\sqrt{2}}{2}$
\choice $3\sqrt{2}$
\end{choices}

---
tipo: Múltipla Escolha
dificuldade: Média
nivel: médio
disciplina: Matemática
assunto: Função Afim
gabarito: C
resposta:
tags: [inequação, custo]
fonte: concurso
concurso:
banca: UFF
ano: 2000
numero: 10
cargo:
prova:
---
\question
As empresas ALFA e BETA alugam televisores do mesmo tipo. A empresa ALFA cobra R$ 35,00 fixos pelos primeiros 30 dias de uso e R$ 1,00 por dia extra. A empresa BETA cobra R$ 15,00 pelos primeiros 20 dias de uso e R$ 1,50 por dia extra. Após $n$ dias o valor cobrado pela empresa BETA passa a ser maior do que o cobrado pela empresa ALFA. O valor de $n$ é:
\begin{choices}
\choice 25
\choice 35
\CorrectChoice 40
\choice 45
\choice 50
\end{choices}

---
tipo: Múltipla Escolha
dificuldade: Fácil
nivel: médio
disciplina: Matemática
assunto: Probabilidade
gabarito: C
resposta:
tags: [análise combinatória]
fonte: concurso
concurso:
banca: UFF
ano: 2000
numero: 11
cargo:
prova:
---
\question
Em uma bandeja há dez pastéis dos quais três são de carne, três de queijo e quatro de camarão. Se Fabiana retirar, aleatoriamente e sem reposição, dois pastéis desta bandeja, a probabilidade de os dois pastéis retirados serem de camarão é:
\begin{choices}
\choice $\frac{3}{25}$
\choice $\frac{4}{25}$
\CorrectChoice $\frac{2}{15}$
\choice $\frac{2}{5}$
\choice $\frac{4}{5}$
\end{choices}

---
tipo: Múltipla Escolha
dificuldade: Difícil
nivel: superior
disciplina: Matemática
assunto: Cônicas
gabarito: E
resposta:
tags: [equação geral]
fonte: concurso
concurso:
banca: UFF
ano: 2000
numero: 12
cargo:
prova:
---
\question
Considere a equação

$(m + n - 1)x^2 + (m - n + 1)y^2 + 2x + 2y - 2 = 0$.

Pode-se afirmar que:
\begin{choices}
\choice Se $m = 0$ e $n = 2$ então a equação representa uma elipse.
\choice Se $m = n = 0$ então a equação representa uma reta.
\choice Se $m = 0$ e $n = 1$ então a equação representa uma parábola.
\choice Se $m = 1$ e $n = 2$ então a equação representa uma hipérbole.
\CorrectChoice Se $m = n = 1$ então a equação representa uma circunferência.
\end{choices}

---
tipo: Múltipla Escolha
dificuldade: Média
nivel: médio
disciplina: Matemática
assunto: Função Logarítmica
gabarito: B
resposta:
tags: [gráfico, distância, figura]
fonte: concurso
concurso:
banca: UFF
ano: 2000
numero: 13
cargo:
prova:
---
\question
A figura representa o gráfico da função $f$ definida por $f(x)=\log_2 x$.

\includegraphics{figura-5}

A medida do segmento $PQ$ é igual a:
\begin{choices}
\choice $\sqrt{6}$
\CorrectChoice $\sqrt{5}$
\choice $\log_2 5$
\choice 2
\choice $\log 2$
\end{choices}

---
tipo: Múltipla Escolha
dificuldade: Média
nivel: médio
disciplina: Matemática
assunto: Juros Compostos
gabarito: A
resposta:
tags: [porcentagem, reajuste]
fonte: concurso
concurso:
banca: UFF
ano: 2000
numero: 14
cargo:
prova:
---
\question
A empresa ACME concedeu a seus funcionários mensalmente, durante dois meses, um reajuste fixo de $x\%$ ao mês. Se ao final desses dois meses o reajuste acumulado foi de 21\%, o valor de $x$ é:
\begin{choices}
\CorrectChoice 10
\choice 10,5
\choice 11
\choice 11,5
\choice 21
\end{choices}

---
tipo: Múltipla Escolha
dificuldade: Média
nivel: médio
disciplina: Matemática
assunto: Sistemas de Equações
gabarito: A
resposta:
tags: [contagem]
fonte: concurso
concurso:
banca: UFF
ano: 2000
numero: 15
cargo:
prova:
---
\question
Cada filha de Luiz Antônio tem o número de irmãs igual à quarta parte do número de irmãos. Cada filho de Luiz Antônio tem o número de irmãos igual ao triplo do número de irmãs.

O total de filhas de Luiz Antônio é:
\begin{choices}
\CorrectChoice 5
\choice 6
\choice 11
\choice 16
\choice 21
\end{choices}

---
tipo: Múltipla Escolha
dificuldade: Média
nivel: médio
disciplina: Matemática
assunto: Progressão Geométrica
gabarito: C
resposta:
tags: [área, razão, figura]
fonte: concurso
concurso:
banca: UFF
ano: 2000
numero: 16
cargo:
prova:
---
\question
Os retângulos $R_1$, $R_2$ e $R_3$, representados na figura, são congruentes e estão divididos em regiões de mesma área.

\includegraphics{figura-6}

Ao se calcular o quociente entre a área da região pintada e a área total de cada um dos retângulos $R_1$, $R_2$ e $R_3$, verifica-se que os valores obtidos formam uma progressão geométrica decrescente de três termos.

A razão dessa progressão geométrica é:
\begin{choices}
\choice $\frac{1}{8}$
\choice $\frac{1}{4}$
\CorrectChoice $\frac{1}{2}$
\choice 2
\choice 4
\end{choices}

---
tipo: Múltipla Escolha
dificuldade: Média
nivel: médio
disciplina: Matemática
assunto: Funções
gabarito: E
resposta:
tags: [gráfico, figura]
fonte: concurso
concurso:
banca: UFF
ano: 2000
numero: 17
cargo:
prova:
---
\question
O gráfico da função $f$ está representado na figura:

\includegraphics{figura-7}

Sobre a função $f$ é FALSO afirmar que:
\begin{choices}
\choice $f(1) + f(2) = f(3)$
\choice $f(2) = f(7)$
\choice $f(3) = 3f(1)$
\choice $f(4) - f(3) = f(1)$
\CorrectChoice $f(2) + f(3) = f(5)$
\end{choices}

---
tipo: Discursiva
dificuldade: Difícil
nivel: superior
disciplina: Matemática
assunto: Polinômios
gabarito:
resposta: "a) z = 0 ou z = 2i ou z = -2i; b) k = -3/2 e r(x) = 19x/2 + 1/2."
tags: [números complexos, divisão de polinômios]
fonte: concurso
concurso:
banca: UFF
ano: 2000
numero: 18
cargo:
prova:
---
\question
Considere os polinômios

$p(x) = 2x^3 + 2x^2 + 7x - 1$ e

$q(x) = 2x^2 - x - 1$.

Calcule:

\begin{parts}
\part os valores do número complexo $z$ tais que $p(z)=q(z)$;
\part o número real $k$ e o polinômio do primeiro grau $r(x)$ tais que $p(x)=(x-k)q(x)+r(x)$.
\end{parts}

---
tipo: Discursiva
dificuldade: Média
nivel: médio
disciplina: Matemática
assunto: Geometria Analítica
gabarito:
resposta: "A = (-5, 0), B = (0, -2), C = (-1, -4)"
tags: [reta, coordenadas]
fonte: concurso
concurso:
banca: UFF
ano: 2000
numero: 19
cargo:
prova:
---
\question
Com relação ao triângulo $ABC$ sabe-se que:

- o ponto $A$ pertence ao eixo das abscissas;
- o ponto $B$ pertence ao eixo das ordenadas;
- a equação da reta que contém os pontos $A$ e $C$ é $x + y + 5 = 0$;
- a equação da reta que contém os pontos $B$ e $C$ é $2x - y - 2 = 0$.

Determine as coordenadas dos pontos $A$, $B$ e $C$.
